
%%%%%%%%%%%%%%%%%%%%%%%%%%%%%%%%%%%%%%%%
%           main body
%%%%%%%%%%%%%%%%%%%%%%%%%%%%%%%%%%%%%%%%

%-----------------------------------
%	TITLE PAGE
%-----------------------------------

\title[实变函数\\
复习]{\texorpdfstring{\Huge \bfseries 实变函数复习}{实变函数复习}} % The short title appears at the bottom of every slide, the full title is only on the title page
% \author[\S 11.3\\
% 多元连续函数性质] % Your institution as it will appear on the bottom of every slide, maybe shorthand to save space
% {\Large \bfseries \S 11.3 多元连续函数的性质}
% \author{Singular Value Decomposition} % Your name
% \date{\today} % Date, can be changed to a custom date
\date{}

%------------------------------------------------

\begin{document}

%------------------------------------------------

\begin{frame}

\titlepage % Print the title page as the first slide

\end{frame}

%------------------------------------------------

\begin{frame}
\frametitle{内容提要} % Table of contents slide, comment this block out to remove it
\tableofcontents % Throughout your presentation, if you choose to use \section{} and \subsection{} commands, these will automatically be printed on this slide as an overview of your presentation
\end{frame}

%-----------------------------------
%	PRESENTATION SLIDES
%-----------------------------------

\section[集合论]{集合论}

%------------------------------------------------

\begin{frame}{集合运算}

集合运算主要包括~\alert{交、并、差、补}

\pause
\vspace{0.6em}

\begin{columns}
\begin{column}[T]{0.55\textwidth}
\begin{block}{运算法则 ~ \reflectbox{$\longsquiggly$} ~ De Morgan 律}
\begin{itemize}
\item $E \cap \left( \bigcup\limits_{\alpha \in I} A_{\alpha} \right) = \bigcup\limits_{\alpha \in I} \left( E \cap A_{\alpha} \right)$
\item $E \cup \left( \bigcap\limits_{\alpha \in I} A_{\alpha} \right) = \bigcap\limits_{\alpha \in I} \left( E \cup A_{\alpha} \right)$
\item $\mathscr{C} \left( \bigcup\limits_{\alpha \in I} A_{\alpha} \right) = \bigcap\limits_{\alpha \in I} \mathscr{C} A_{\alpha}$
\item $\mathscr{C} \left( \bigcap\limits_{\alpha \in I} A_{\alpha} \right) = \bigcup\limits_{\alpha \in I} \mathscr{C} A_{\alpha}$
\end{itemize}
\end{block}
\end{column}
\pause
\begin{column}[T]{0.4\textwidth}
特别注意: 图解法
\end{column}
\end{columns}

\end{frame}

%------------------------------------------------

\begin{frame}{集合的对等、势}

\end{frame}

%------------------------------------------------

\begin{frame}{开、闭集相关的概念与性质}

\end{frame}

%------------------------------------------------

\begin{frame}{Cantor 三分集的构造}

{\smaller[1]
\begin{align*}
F_1 & = F_{11} \cup F_{12} = \left[ 0, \dfrac{1}{3} \right] \cup \left[ \dfrac{2}{3}, 1 \right], \\
I_1 & = I_{11} = \left( \dfrac{1}{3}, \dfrac{2}{3} \right), \\
F_2 & = F_{21} \cup F_{22} \cup F_{23} \cup F_{24} = \left[ 0, \dfrac{1}{9} \right] \cup \left[ \dfrac{2}{9}, \dfrac{1}{3} \right] \cup \left[ \dfrac{2}{3}, \dfrac{7}{9} \right] \cup \left[ \dfrac{8}{9}, 1 \right], \\
I_2 & = I_{21} \cup I_{22} = \left( \dfrac{1}{9}, \dfrac{2}{9} \right) \cup \left( \dfrac{7}{9}, \dfrac{8}{9} \right), \\
& \vdots \\
F_n & = F_{n1} \cup F_{n2} \cup \cdots \cup F_{n2^{n}}, \\
I_n & = I_{n1} \cup I_{n2} \cup \cdots \cup I_{n2^{n-1}}, \\
& \vdots \\
G_0 & = \bigcup_{n=1}^{\infty} I_n, \\
P_0 & = \mathscr{C} G_0 = \bigcap_{n=1}^{\infty} F_n \longleftarrow \text{(康托三分集)}.
\end{align*}
}

\end{frame}

%------------------------------------------------

\begin{frame}{Cantor 三分集的性质}

\end{frame}

%------------------------------------------------

\section[测度论]{测度论 (可测集、可测函数)}

%------------------------------------------------

\begin{frame}{开、闭集测度}



\end{frame}

%------------------------------------------------

\begin{frame}{内、外测度}



\end{frame}

%------------------------------------------------

\begin{frame}{可测集}



\end{frame}

%------------------------------------------------

\begin{frame}{渐缩、渐张可测集列}



\end{frame}

%------------------------------------------------

\begin{frame}{环, $\sigma$ 环, $\sigma$ 代数及其上的测度}



\end{frame}

%------------------------------------------------

\begin{frame}{可测函数}



\end{frame}

%------------------------------------------------

\begin{frame}{可测函数的运算}



\end{frame}

%------------------------------------------------

\begin{frame}{Egorov 定理}



\end{frame}

%------------------------------------------------

\begin{frame}{Riesz 定理}



\end{frame}

%------------------------------------------------

\begin{frame}{Luzin 定理}



\end{frame}

%------------------------------------------------

\begin{frame}{复合函数可测性}



\end{frame}

%------------------------------------------------

\section[积分论]{积分论}

%------------------------------------------------

\begin{frame}{勒贝格积分的构造}



\end{frame}

%------------------------------------------------

\begin{frame}{单个积分的性质}



\end{frame}

%------------------------------------------------

\begin{frame}{积分序列的性质 | {\smaller[1] Levi 定理}}



\end{frame}

%------------------------------------------------

\begin{frame}{积分序列的性质 | {\smaller[1] Fatou 引理}}



\end{frame}

%------------------------------------------------

\begin{frame}{积分序列的性质 | {\smaller[1] Lebesgue 控制收敛定理}}



\end{frame}

%------------------------------------------------

\begin{frame}{积分序列的性质 | {\smaller[1] Lebesgue-Vitali 依测度型控制收敛定理}}



\end{frame}

%------------------------------------------------

\begin{frame}{Riemann 积分与 Lebesgue 积分比较}



\end{frame}

%------------------------------------------------

\begin{frame}{多重积分与累次积分 (Fubini 定理)}



\end{frame}

%------------------------------------------------

\begin{frame}{(变上限) 积分与微分}

% Lebesgue 积分与微分的关系可用如下交换图表总结:

\begin{adjustbox}{width=1.0\textwidth}
\begin{tikzpicture}[
% scale=\tikzscale,
transform shape,
sarr/.style={-Classical TikZ Rightarrow[]},
darr/.style={-Classical TikZ Rightarrow[sep] Classical TikZ Rightarrow[]},
dsarr/.style={sarr, dashed},
ddarr/.style={darr, dashed}
]

\node (L0) at (0, 0) {$L_0$};
\node (L) at (2, 0) {$L_{[a, b]}$};
\node (B) at (7, 0) {$C([a, b])$};
\node (BV) at (7, -2) {$BV([a, b]) \cap C([a, b])$};
\node (BV2) at (11, -2) {$BV([a, b])$};
\node (AC) at (7, -4) {$AC_0([a, b])$};
\node (p1) at (2, -4) {$L_{[a, b]} / L_0$};
\node (L_again) at (11, -4) {$L_{[a, b]}$};
\node (p2) at (14, -4) {$L_{[a, b]} / L_0$};
\node (AC2) at (17, -4) {$AC_0([a, b])$};
\node (ker1) at (15, -2) {$\ker \left( {\textstyle\int}_{[a,x]}~\circ~\operatorname{pr}~\circ~\widetilde{\mathrm{d}} \right)$};
\node (ker2) at (11, 0) {$\ker \left( {\textstyle\int}_{[a,x]}~\circ~\operatorname{pr}~\circ~\widetilde{\mathrm{d}}~\circ~\iota \right)$};

\draw[sarr] ([xshift=0.5ex,yshift=1ex] L0.east) arc (90:270:0.5ex) -- (L);
\draw[sarr] (L) -- (B) node[midway,above] {${\textstyle\int}_{[a, x]}$};
\draw[sarr] ([xshift=-1ex,yshift=0.5ex] BV.north) arc (180:360:0.5ex) -- (B);
\draw[dsarr] (L) -- (BV) node[pos=0.5,above] {${\textstyle\int}_{[a, x]}$} node[pos=0.9,above] {\Circled{1}};
\draw[sarr] ([xshift=0.5ex,yshift=1ex] BV.east) arc (90:270:0.5ex) -- (BV2) node[midway,above] {$\iota$};
\draw[sarr] ([xshift=-1ex,yshift=0.5ex] AC.north) arc (180:360:0.5ex) -- (BV) node[midway,right] {\Circled{2}};
\draw[ddarr] (L) -- ([xshift=-4ex] AC.north) node[midway,left] {${\textstyle\int}_{[a, x]}$}  node[pos=0.9,above] {\Circled{3}};
\draw[darr] (L) -- (p1) node[midway,left] {$\operatorname{pr}$};
\draw[dsarr] (p1) -- (AC) node[midway,above] {${\textstyle\int}_{[a, x]}$} node[midway,below] {$\sim$};
\draw[dsarr] (AC) -- (L_again) node[midway,above] {$\widetilde{~ \mathrm{d}}$};
\draw[dsarr] (BV2) -- (L_again) node[midway,left] {$\widetilde{~ \mathrm{d}}$} node[midway,right] {\Circled{4}};
\draw[darr] (L_again) -- (p2) node[midway,above] {$\operatorname{pr}$};
\draw[sarr] (p2) -- (AC2) node[midway,above] {${\textstyle\int}_{[a,x]}$};
\draw[dsarr, bend right = 15] (p1) to node[midway,below] {$\operatorname{pr}~\circ~\widetilde{~ \mathrm{d}}~\circ~{\textstyle\int}_{[a, x]} = \operatorname{id}$} node[midway,above] {\Circled{5}} (p2);
\draw[sarr] ([xshift=-0.5ex,yshift=1ex] ker1.west) arc (90:-90:0.5ex) -- (BV2);
\draw[sarr] ([xshift=-0.5ex,yshift=1ex] ker2.south west) arc (120:-60:0.5ex) -- (BV);

\end{tikzpicture}

\end{adjustbox}

\vspace{-2em}

\begin{align*}
L_{[a, b]} & = [a, b] \text{ 区间上的 Lebesgue 可积函数}, \\
L_0 & = \{ f \in L_{[a, b]} ~:~ f \sim 0\}, \\
C([a, b]) & = [a, b] \text{ 区间上的连续函数}, \\
AB([a, b]) & = [a, b] \text{ 区间上几乎处处有限的函数}, \\
BV([a, b]) & = [a, b] \text{ 区间上有界变差函数}, \\
AC([a, b]) & = [a, b] \text{ 区间上绝对连续函数}.
\end{align*}

\end{frame}

%------------------------------------------------

\begin{frame}{(变上限) 积分与微分}

\end{frame}

%------------------------------------------------

\section[$L^p$ 空间]{$L^p$ 空间}

%------------------------------------------------

\begin{frame}{$L^p$ 空间定义与基本性质}



\end{frame}

%------------------------------------------------

\begin{frame}{Hölder 不等式}



\end{frame}

%------------------------------------------------

\begin{frame}{Minkowski 不等式}



\end{frame}

%------------------------------------------------

\begin{frame}{强收敛、弱收敛}



\end{frame}

%------------------------------------------------

\begin{frame}{各种意义下收敛的关系}

\begin{adjustbox}{width=1.0\textwidth}
\begin{tikzpicture}[
% scale=1,
draw=black,
text=black,
transform shape
]

\node (uniform) at (-12, -4) {$\text{({\color{cyan}近})一致收敛}$};
\node (ae) at (-7,0) {$\text{({\color{magenta}子列})几乎处处收敛}$};
\node (measure) at (5, 0) {$\text{依测度收敛}$};
\node (norm) at (-3, -4) {$\text{强收敛（依范数收敛）}$};
\node (weak) at (-3, -8) {$\text{弱收敛}$};

\draw[arrows={-{Stealth[length=3mm, width=2mm]}}, transform canvas={xshift=0.5em}] (uniform) -- (ae) node[sloped, anchor=center, midway, below] {$\checkmark$};
\draw[arrows={-{Stealth[length=3mm, width=2mm]}}, transform canvas={xshift=-0.5em}] (ae) -- (uniform) node[sloped, anchor=center, midway, above] {{\color{red}$\boldsymbol{\times}$}, ~~ {\color{cyan} Egorov ($m E < \infty$)}};

\draw[arrows={-{Stealth[length=3mm, width=2mm]}}, transform canvas={yshift=0.3em}] (ae) -- (measure) node[sloped, anchor=center, midway, above] {{\color{red}$\boldsymbol{\times}$}, ~~ $m E < \infty$};
\draw[arrows={-{Stealth[length=3mm, width=2mm]}}, transform canvas={yshift=-0.3em}] (measure) -- (ae) node[sloped, anchor=center, midway, below] {{\color{red}$\boldsymbol{\times}$}, ~~ {\color{magenta} Riesz ($m E < \infty$)}};

\draw[arrows={-{Stealth[length=3mm, width=2mm]}}, transform canvas={xshift=0.5em}] (norm) -- (ae) node[sloped, anchor=center, midway, above] {\color{red}$\boldsymbol{\times}$};
\draw[arrows={-{Stealth[length=3mm, width=2mm]}}, transform canvas={xshift=-0.5em}] (ae) -- (norm) node[sloped, anchor=center, midway, below] {{\color{red}$\boldsymbol{\times}$}, ~ $\lVert f_n \rVert_p \to \lVert f \rVert_p$};

\draw[arrows={-{Stealth[length=3mm, width=2mm]}}, transform canvas={xshift=-0.8em}] (norm) -- (measure) node[sloped, anchor=center, midway, above] {$\checkmark$};
\draw[arrows={-{Stealth[length=3mm, width=2mm]}}, transform canvas={xshift=0.8em}] (measure) -- (norm) node[sloped, anchor=center, midway, below] {{\color{red}$\boldsymbol{\times}$}, ~~ $\text{\smaller 等度绝对连续积分}$ ($m E < \infty$)};

\draw[arrows={-{Stealth[length=3mm, width=2mm]}}, transform canvas={xshift=-0.3em}] (norm) -- (weak) node[midway, left] {$\checkmark$};
\draw[arrows={-{Stealth[length=3mm, width=2mm]}}, transform canvas={xshift=0.3em}] (weak) -- (norm) node[midway, right] {\color{red}$\boldsymbol{\times}$};

\draw[arrows={-{Stealth[length=3mm, width=2mm]}}, transform canvas={yshift=-0.3em}] (uniform) -- (weak) node[sloped, anchor=center, midway, below] {{\color{red}$\boldsymbol{\times}$}, ~~ $m E < \infty$};
\draw[arrows={-{Stealth[length=3mm, width=2mm]}}, transform canvas={yshift=0.3em}] (weak) -- (uniform) node[sloped, anchor=center, midway, above] {\color{red}$\boldsymbol{\times}$};

\draw[arrows={-{Stealth[length=3mm, width=2mm]}}, transform canvas={yshift=0.3em}] (uniform) -- (norm) node[sloped, anchor=center, midway, above] {{\color{red}$\boldsymbol{\times}$}, ~~ $m E < \infty$};
\draw[arrows={-{Stealth[length=3mm, width=2mm]}}, transform canvas={yshift=-0.3em}] (norm) -- (uniform) node[sloped, anchor=center, midway, below] {{\color{red}$\boldsymbol{\times}$}};

\draw[arrows={-{Stealth[length=3mm, width=2mm]}}, transform canvas={xshift=0.7em}] (measure) -- (weak) node[sloped, anchor=center, midway, above] {{\color{red}$\boldsymbol{\times}$}};
\draw[arrows={-{Stealth[length=3mm, width=2mm]}}, transform canvas={xshift=1.6em}] (weak) -- (measure) node[sloped, anchor=center, midway, below] {{\color{red}$\boldsymbol{\times}$}};

\end{tikzpicture}

\end{adjustbox}

$A \xrightarrow{{\color{red} \boldsymbol{\times}}} B$ 表示 $A$ 不蕴含 (不能推出) $B,$ 如果 $\color{red} \times$ 之后有注释, 则表示在某些条件下成立, 相应的结论也可能会弱一些 (对应括号中同颜色的文字). $A \xrightarrow{\checkmark} B$ 表示 $A$ 蕴含 $B.$ 还有要注意的是, 讨论蕴含关系时, 都是在某个空间中进行讨论的, 例如, $L^p$ 空间等.

\end{frame}

%------------------------------------------------

\begin{frame}{各种意义下收敛的关系 | {\smaller[1] 一些典型的反例}}

\begin{itemize}[<+->]
\item 几乎处处收敛但不依测度收敛的例子 (注意, 必须有 $m E = \infty$):
$$f_n(x) = \chi_{[n, n + 1]}(x), \quad n \in \mathbb{N}.$$

\item 依测度收敛但不几乎处处收敛的例子: $E = [0, 1],$
$$f_n(x) = \chi_{\left[ \frac{i}{2^k}, \frac{i + 1}{2^k} \right]}(x), \quad n = 2^k + i, \quad 0 \leqslant i < 2^k.$$
以上也是强收敛但不几乎处处收敛的例子.

\item 依测度收敛但不强收敛的例子 (稍微修改上例):
$$f_n(x) = {\color{red} 2^k} \chi_{\left[ \frac{i}{2^k}, \frac{i + 1}{2^k} \right]}(x), \quad n = 2^k + i, \quad 0 \leqslant i < 2^k.$$
依测度收敛 (到 $0$) 但不弱收敛的例子 (取 $g(x) = 1$).
\end{itemize}

\end{frame}

%------------------------------------------------

\begin{frame}{各种意义下收敛的关系 | {\smaller[1] 一些典型的反例}}

\begin{itemize}[<+->]
\item 一致收敛但不弱收敛 (也不强收敛) 的例子 (必须有 $m E = \infty$):
$$f_n(x) = \dfrac{1}{n} \chi_{[1, e^n]}(x).$$
容易验证 $f_n(x)$ 一致收敛到 $0.$ 设 $p, q > 1$ 为共轭数, 则 $f_n \in L^p(\mathbb{R}).$ 取 $\displaystyle g(x) = \dfrac{\chi_{[1, +\infty)}(x)}{x},$ 则 $g \in L^q(\mathbb{R}),$ 但是 $\displaystyle \int_{\mathbb{R}} f_n(x) g(x) ~ \mathrm{d} x = \dfrac{1}{n} \int_{1}^{e^n} \dfrac{1}{x} ~ \mathrm{d} x = 1.$

\item 弱收敛但不一致收敛的例子: $E = [0, 1),$
$$f_n(x) = x^n, \quad n \in \mathbb{N}.$$

\item 弱收敛但不强收敛的例子: $E = (0, \pi),$
$$f_n(x) = \sin(nx), \quad n \in \mathbb{N}.$$
\end{itemize}

\end{frame}

%------------------------------------------------

\begin{frame}{依 $L^1$ 意义傅里叶变换}

\end{frame}

%------------------------------------------------

\begin{frame}{依 $L^2$ 意义傅里叶变换}

\end{frame}

%------------------------------------------------

\end{document}
